% Volume II: Causal Explosion
% Part IV. Time, Delay, and the Maturation of Evidence
% Chapter 18: Millions of Hours
% Spine: proof-spines/ch018-spine.md

\chapter{Millions of Hours}
\label{ch:ch018}

Modern inquiries often confront \textbf{millions of hours} of record rather than a scarcity of traces. If a public event produces \(V\) hours of video and reviewers can process footage at an effective rate \(r\), then naïve serial review requires
\[
V/r
\]
\ToyModel\ This serial-review quotient isolates the throughput bottleneck before coordination costs are added.
hours of attention. Parallel review lowers wall-clock time, but only by introducing coordination burdens, duplication, inconsistent labeling, and escalation costs. This is \textbf{capacity delay} in its purest form: the institution withholds no evidence and gains nothing from waiting, it simply does not yet have enough throughput to convert what it holds into judgment.

The chapter therefore models evidentiary review as a queue. Evidence items arrive at rate \(\lambda\), reviewers process them at rate \(\mu\), and only a fraction \(p\) are escalated for higher scrutiny. When arrival outstrips effective review capacity, backlog grows even under comprehensive recording. The lesson is that duration can be probative while still being operationally hard to harvest: a society may capture immense evidence and yet fail to convert it into timely judgment.
